\chapter{Results}

The values given in this chapter are the means of three different leaves sampled from each pot, except for the fresh, turgid and dry weights, which were measured on one leaf per pot. Each combination of species, growing medium and salinity level is represented by a single pot, so the differences between treatments are described and the effects of biochar and salinity were tested within each species on the pot means (Section~\ref{sec:anova}). Percentage differences are calculated from the means in Tables~\ref{tab:growth} and~\ref{tab:physiology}.

\section{Growth responses of mint species to salt stress}

Salinity reduced the growth of all three taxa (Table~\ref{tab:growth}). In unamended peat, plant height at 50~mM NaCl was 31\% lower than at 0~mM in \textit{M. piperita}, 44\% lower in \textit{M. crispata} and 51\% lower in \textit{M. spicata} `Moroccan'. Leaf length and leaf width followed the same pattern: in unamended peat, leaf length at 50~mM was 21\%, 55\% and 33\% lower than at 0~mM, and leaf width 19\%, 44\% and 24\% lower, in the same order of taxa. The reduction was already visible at 25~mM, where the values were in most cases intermediate between those of the 0 and 50~mM treatments.

\textit{M. piperita} was the tallest taxon (36.1 to 64.4~cm), followed by \textit{M. crispata} (24.0 to 53.2~cm), while \textit{M. spicata} `Moroccan' was clearly the shortest (9.1 to 22.8~cm).

\subsection{\textit{Mentha piperita}}

In unamended peat, \textit{M. piperita} reached 52.11~cm without salt, 44.17~cm at 25~mM and 36.12~cm at 50~mM, corresponding to losses of 15.2 and 30.7\% respectively. The decline between the two salinity levels was therefore steeper than that between the control and 25~mM, which suggests that the lower concentration was tolerated comparatively well. Leaf dimensions were less affected than height: leaf length fell from 5.77 to 5.46~cm at 25~mM and to 4.55~cm at 50~mM, losses of 5.4 and 21.1\%, and leaf width from 3.24 to 3.14 and 2.62~cm, losses of 3.1 and 19.1\%. At 25~mM the leaves of this taxon were thus almost unchanged while height had already fallen by about a sixth, indicating that stem extension responded to the salt before leaf expansion did.

\subsection{\textit{Mentha crispata}}

\textit{M. crispata} was more strongly affected. Height in unamended peat fell from 43.00~cm to 35.90~cm at 25~mM and to 24.00~cm at 50~mM, losses of 16.5 and 44.2\%. Unlike in \textit{M. piperita}, the leaves of this taxon responded as strongly as the stems: leaf length fell from 4.65 to 3.89 and then 2.11~cm, a loss of 54.6\% at the highest salinity, and leaf width from 2.50 to 2.20 and 1.40~cm, a loss of 44.0\%. The leaf length of this taxon showed the largest proportional reduction recorded for any growth parameter in the experiment. The plants at 50~mM in unamended peat were therefore not only shorter but carried markedly smaller leaves, so that the reduction in total leaf area would have been greater than the reduction in height alone implies.

\subsection{\textit{Mentha spicata} `Moroccan'}

\textit{M. spicata} `Moroccan' was the smallest taxon throughout and the most sensitive in terms of height. In unamended peat it fell from 18.50~cm to 13.00~cm at 25~mM, a loss of 29.7\%, which is roughly twice the proportional loss shown by the other two taxa at the same concentration, and to 9.10~cm at 50~mM, a loss of 50.8\%. This taxon was thus the only one to lose close to a third of its height at the intermediate salinity. Its leaves were comparatively less affected, leaf length falling by 12.8 and 33.4\% and leaf width by 2.9 and 23.8\% at 25 and 50~mM. The pattern of this taxon therefore resembles that of \textit{M. piperita}, with height responding before leaf size, but with the height response much more pronounced.

\subsection{Comparison between the taxa}

Ranking the taxa by their proportional loss of height at 50~mM in unamended peat gives \textit{M. spicata} `Moroccan' (50.8\%), \textit{M. crispata} (44.2\%) and \textit{M. piperita} (30.7\%), so that the tallest taxon was also the least affected in relative terms. The ranking by leaf length is different, \textit{M. crispata} losing most (54.6\%), followed by `Moroccan' (33.4\%) and \textit{M. piperita} (21.1\%). The taxa therefore did not differ simply in how much salt they tolerated, but in which organ responded first: height was the more sensitive measure in \textit{M. piperita} and `Moroccan', whereas in \textit{M. crispata} the leaves were affected at least as strongly as the stems. For all three taxa the values at 25~mM lay between those of the control and the 50~mM treatment for every parameter measured, so no parameter showed a threshold response within the range tested.

\begin{singlespace}
\small
\begin{longtable}{p{3cm}cc S[table-format=2.2,round-mode=places,round-precision=2] S[table-format=1.2,round-mode=places,round-precision=2] S[table-format=1.2,round-mode=places,round-precision=2]}
  \caption[Plant height and leaf size]{Plant height and leaf size of the three mint taxa by biochar rate and salinity level (mean of the leaves sampled from each pot).} \label{tab:growth} \\
  \toprule
  Species & Biochar (g) & NaCl (mM) & {Plant height (cm)} & {Leaf length (cm)} & {Leaf width (cm)} \\
  \midrule
  \endfirsthead
  \toprule
  Species & Biochar (g) & NaCl (mM) & {Plant height (cm)} & {Leaf length (cm)} & {Leaf width (cm)} \\
  \midrule
  \endhead
  \bottomrule
  \endlastfoot
  \textit{M. piperita} & 0 & 0 & 52.11 & 5.77 & 3.24 \\
   & 0 & 25 & 44.17 & 5.46 & 3.14 \\
   & 0 & 50 & 36.12 & 4.55 & 2.62 \\
   & 15 & 0 & 55.78 & 6.41 & 3.36 \\
   & 15 & 25 & 51.11 & 5.82 & 3.26 \\
   & 15 & 50 & 42.77 & 4.9 & 2.92 \\
   & 30 & 0 & 64.41 & 6.15 & 3.62 \\
   & 30 & 25 & 58.14 & 6.43 & 3.5 \\
   & 30 & 50 & 52.81 & 6.01 & 3.32 \\
  \midrule
  \textit{M. crispata} & 0 & 0 & 43 & 4.65 & 2.5 \\
   & 0 & 25 & 35.9 & 3.89 & 2.2 \\
   & 0 & 50 & 24 & 2.11 & 1.4 \\
   & 15 & 0 & 45.97 & 4.7 & 2.54 \\
   & 15 & 25 & 41.45 & 4.5 & 2.54 \\
   & 15 & 50 & 31.65 & 3.52 & 2.09 \\
   & 30 & 0 & 53.24 & 5.36 & 2.64 \\
   & 30 & 25 & 46.62 & 4.96 & 2.73 \\
   & 30 & 50 & 38.47 & 4.44 & 2.39 \\
  \midrule
  \textit{M. spicata} `Moroccan' & 0 & 0 & 18.5 & 4.13 & 2.1 \\
   & 0 & 25 & 13 & 3.6 & 2.04 \\
   & 0 & 50 & 9.1 & 2.75 & 1.6 \\
   & 15 & 0 & 20.69 & 4.59 & 2.19 \\
   & 15 & 25 & 18 & 4.18 & 2.1 \\
   & 15 & 50 & 13.28 & 3.2 & 1.83 \\
   & 30 & 0 & 22.81 & 4.97 & 2.4 \\
   & 30 & 25 & 20.26 & 4.52 & 2.19 \\
   & 30 & 50 & 16.33 & 3.89 & 1.93 \\
\end{longtable}
\end{singlespace}

\section{Effects of growing media on plant growth under salt stress}

The biochar rates of 15~g and 30~g correspond to 3.75\% and 7.5\% of the 400~g of substrate in each pot. Plants grown with biochar were taller and had larger leaves than plants in unamended peat at every salinity level (Table~\ref{tab:growth}). Without salt (0~mM), plant height was 7\%, 7\% and 12\% higher with 15~g of biochar, and 24\%, 24\% and 23\% higher with 30~g, than in unamended peat for \textit{M. piperita}, \textit{M. crispata} and \textit{M. spicata} `Moroccan', respectively. The difference between the biochar rates and unamended peat was larger under salt stress: at 50~mM, plant height with 30~g of biochar was 46\%, 60\% and 79\% higher than in unamended peat for the three taxa, respectively, and with 15~g it was 18\%, 32\% and 46\% higher.

Biochar also reduced the loss of growth caused by salinity. The decrease in plant height between 0 and 50~mM was 31\%, 23\% and 18\% in \textit{M. piperita} with 0, 15 and 30~g of biochar, respectively; 44\%, 31\% and 28\% in \textit{M. crispata}; and 51\%, 36\% and 28\% in \textit{M. spicata} `Moroccan'. The same trend was found for leaf length and width, with two exceptions: the decrease in leaf length of \textit{M. piperita} was larger with 15~g of biochar (24\%) than in unamended peat (21\%), and the decrease in leaf width of \textit{M. spicata} `Moroccan' was larger with 30~g (20\%) than with 15~g (16\%), although still smaller than in unamended peat (24\%). For example, in \textit{M. crispata} the leaf length at 50~mM was 55\% lower than at 0~mM in unamended peat, but only 25\% and 17\% lower with 15 and 30~g of biochar.

\subsection{Dependence of the biochar effect on salinity}

The advantage conferred by biochar was not constant across the salinity levels but increased as the salt concentration rose. Expressing the height of the amended plants as a percentage of the height of the unamended plants at the same salinity makes this apparent. Without salt, the 30~g rate gave plants 23.6\% taller in \textit{M. piperita}, 23.8\% taller in \textit{M. crispata} and 23.3\% taller in \textit{M. spicata} `Moroccan', so that the three taxa benefited to almost exactly the same degree. At 25~mM the same comparison gave advantages of 31.6, 29.9 and 55.8\%, and at 50~mM of 46.2, 60.3 and 79.5\%.

The 15~g rate behaved in the same way at a lower level. Its advantage over unamended peat rose from 7.0, 6.9 and 11.8\% without salt to 18.4, 31.9 and 45.9\% at 50~mM, in the order \textit{M. piperita}, \textit{M. crispata}, \textit{M. spicata} `Moroccan'. In every taxon and at both rates, therefore, the margin between amended and unamended plants widened with each increase in salinity.

Two points follow from this. First, under non-saline conditions the three taxa responded almost identically to the higher biochar rate, differing by less than half a percentage point, whereas under salt their responses diverged widely. The differences between taxa in their reaction to biochar thus emerged only when the plants were stressed. Second, the largest advantage recorded, 79.5\% in \textit{M. spicata} `Moroccan' at 50~mM, occurred in the taxon that performed worst in unamended peat under the same conditions; the amended plants of that taxon at 50~mM reached 16.33~cm against 9.10~cm in unamended peat, so that biochar recovered a substantial part of the growth that the salt had removed.

\section{Physiological responses of mint species to salt stress}

The Fv/Fm ratio ranged from 0.787 to 0.865 (Table~\ref{tab:physiology}). It was slightly lower at 50~mM than at 0~mM in all taxa and growing media, by 3 to 7\%, and the biochar rate changed it by at most 4\% at a given salinity level.

The size of this change deserves comment. In unamended peat the ratio fell from 0.863 to 0.820 in \textit{M. piperita}, from 0.841 to 0.815 in \textit{M. crispata} and from 0.850 to 0.787 in \textit{M. spicata} `Moroccan', reductions of 5.0, 3.1 and 7.4\%. Values between 0.78 and 0.86 lie within or only a little below the range usually reported for unstressed leaves, so the photosystem remained largely intact even at the highest salinity, while plant height over the same interval fell by between 30 and 51\%. Growth was therefore affected far more than the efficiency of photosystem~II, which indicates that the reduction in growth was not primarily caused by damage to the photosynthetic apparatus. The lowest value of the whole experiment, 0.787, was recorded in \textit{M. spicata} `Moroccan' in unamended peat at 50~mM, the single combination in which the plants also lost most height, so the two measures agree in identifying that treatment as the most severe.

The chlorophyll content decreased more strongly with salinity. In unamended peat, the chlorophyll content at 50~mM was 29\% lower than at 0~mM in \textit{M. piperita}, 51\% lower in \textit{M. crispata} and 59\% lower in \textit{M. spicata} `Moroccan'. With 30~g of biochar the decreases were 11\%, 21\% and 28\%, respectively. The intermediate rate gave intermediate values, the decreases with 15~g being 18.3, 22.8 and 39.5\% in the same order, so that chlorophyll followed the same graded response to the biochar rate as plant height did. At 50~mM, the chlorophyll content with 30~g of biochar was 33\%, 83\% and 109\% higher than in unamended peat for the three taxa, respectively. One exception to the general trend was \textit{M. crispata} in unamended peat, where the chlorophyll content at 25~mM (30.05) was higher than at 0~mM (21.67). This single value is out of keeping with every other series in the table: in the eight other combinations of taxon and growing medium the chlorophyll content fell at each successive salinity level, and in the amended pots of the same taxon it fell from 22.00 to 20.60 and 16.99 with 15~g and from 24.85 to 23.56 and 19.57 with 30~g. The unamended series for this taxon is also the one that produced the lowest chlorophyll content of the whole experiment at 50~mM, 10.70, so the series both rises and falls more steeply than any other. Whether the value at 25~mM reflects a real response or an error in a single reading cannot be settled from one pot, and the analysis of variance for the chlorophyll content of this taxon, reported in Section~\ref{sec:anova}, returns no significant effect of either factor, which this value is likely to have caused.

Taken together, the two physiological measures behaved differently from one another. The chlorophyll content tracked plant height closely, falling with salinity, rising with the biochar rate, and showing the same graded response to both factors; the proportional decrease in chlorophyll at 50~mM in unamended peat, 29.4\%, 50.6\% and 58.9\% for the three taxa, is close to the corresponding decrease in height, 30.7\%, 44.2\% and 50.8\%. The Fv/Fm ratio, by contrast, moved very little in any treatment. The pigment content of the leaves was therefore reduced considerably while the efficiency of the remaining photosystem~II was largely preserved.

The relative water content (RWC), calculated from the fresh, turgid and dry weights with Equation~\ref{eq:rwc}, ranged from 12.4 to 99.9\% (Table~\ref{tab:physiology}). It was much lower in all pots with biochar (12.4 to 29.3\%) than in the unamended pots of \textit{M. crispata} (73.2 to 83.8\%) and \textit{M. spicata} `Moroccan' (85.5 to 99.9\%), whereas in \textit{M. piperita} the unamended pots also had low values (20.0 to 27.7\%). The RWC did not decrease with salinity: it was higher at 50~mM than at 0~mM in eight of the nine combinations of taxon and growing medium, for example by 38\% in unamended \textit{M. piperita} and by 48\% and 46\% with 15 and 30~g of biochar; the only decrease was in unamended \textit{M. crispata} (5\%). The fresh, turgid and dry weights of the sampled leaf are given in Table~\ref{tab:weights}.

These values are not biologically plausible and are therefore not interpreted here. A relative water content below about 50\% corresponds to severe wilting, and the plants from which the leaves were taken were turgid and alive at sampling, so the figures of 12 to 29\% recorded for the biochar pots cannot represent the true water status of those plants. The direction of the response is also contrary to expectation, the relative water content rising rather than falling with salinity in eight of the nine combinations. The turgid weights recorded for these leaves were four to seven times their fresh weights, which inflates the denominator of Equation~\ref{eq:rwc} and depresses the calculated value; leaves that take up water beyond full turgor, or that retain surface water at weighing, produce exactly this pattern. As the soaking and weighing procedure was the same for every pot, the discrepancy cannot be attributed to differences in method between treatments, and its cause has not been established. The relative water content is reported for completeness but is excluded from the interpretation of the results.

\begin{singlespace}
\small
\begin{longtable}{p{3cm}cc S[table-format=1.3,round-mode=places,round-precision=3] S[table-format=3.2,round-mode=places,round-precision=2] S[table-format=2.1,round-mode=places,round-precision=1]}
  \caption[Fv/Fm, chlorophyll content and RWC]{Fv/Fm ratio, chlorophyll content and relative water content (RWC) of the three mint taxa by biochar rate and salinity level (mean of the leaves sampled from each pot; RWC calculated from the weights in Table~\ref{tab:weights}).} \label{tab:physiology} \\
  \toprule
  Species & Biochar (g) & NaCl (mM) & {Fv/Fm} & {Chlorophyll content ($\mu$mol~m$^{-2}$)} & {RWC (\%)} \\
  \midrule
  \endfirsthead
  \toprule
  Species & Biochar (g) & NaCl (mM) & {Fv/Fm} & {Chlorophyll content ($\mu$mol~m$^{-2}$)} & {RWC (\%)} \\
  \midrule
  \endhead
  \bottomrule
  \endlastfoot
  \textit{M. piperita} & 0 & 0 & 0.863 & 41.18 & 20.0 \\
   & 0 & 25 & 0.843 & 32.89 & 21.0 \\
   & 0 & 50 & 0.82 & 29.06 & 27.7 \\
   & 15 & 0 & 0.86 & 39.05 & 15.1 \\
   & 15 & 25 & 0.843 & 36.26 & 21.4 \\
   & 15 & 50 & 0.822 & 31.91 & 22.4 \\
   & 30 & 0 & 0.865 & 43.6 & 13.4 \\
   & 30 & 25 & 0.865 & 40.54 & 19.6 \\
   & 30 & 50 & 0.83 & 38.69 & 19.6 \\
  \midrule
  \textit{M. crispata} & 0 & 0 & 0.841 & 21.67 & 76.9 \\
   & 0 & 25 & 0.86 & 30.05 & 83.8 \\
   & 0 & 50 & 0.815 & 10.7 & 73.2 \\
   & 15 & 0 & 0.838 & 22 & 24.8 \\
   & 15 & 25 & 0.828 & 20.6 & 26.1 \\
   & 15 & 50 & 0.797 & 16.99 & 29.3 \\
   & 30 & 0 & 0.847 & 24.85 & 24.4 \\
   & 30 & 25 & 0.843 & 23.56 & 24.2 \\
   & 30 & 50 & 0.819 & 19.57 & 27.2 \\
  \midrule
  \textit{M. spicata} `Moroccan' & 0 & 0 & 0.85 & 15.32 & 85.5 \\
   & 0 & 25 & 0.846 & 10.75 & 88.3 \\
   & 0 & 50 & 0.787 & 6.3 & 99.9 \\
   & 15 & 0 & 0.86 & 17.4 & 12.9 \\
   & 15 & 25 & 0.831 & 13.88 & 14.9 \\
   & 15 & 50 & 0.816 & 10.52 & 20.4 \\
   & 30 & 0 & 0.863 & 18.35 & 12.4 \\
   & 30 & 25 & 0.852 & 15.84 & 15.6 \\
   & 30 & 50 & 0.82 & 13.16 & 16.6 \\
\end{longtable}
\end{singlespace}

\begin{singlespace}
\small
\begin{longtable}{p{3cm}cc S[table-format=2.2,round-mode=places,round-precision=2] S[table-format=3.2,round-mode=places,round-precision=2] S[table-format=1.2,round-mode=places,round-precision=2]}
  \caption[Leaf fresh, turgid and dry weight]{Fresh weight (FW), turgid weight (TW) and dry weight (DW) of one leaf of the three mint taxa by biochar rate and salinity level (mg).} \label{tab:weights} \\
  \toprule
  Species & Biochar (g) & NaCl (mM) & {FW (mg)} & {TW (mg)} & {DW (mg)} \\
  \midrule
  \endfirsthead
  \toprule
  Species & Biochar (g) & NaCl (mM) & {FW (mg)} & {TW (mg)} & {DW (mg)} \\
  \midrule
  \endhead
  \bottomrule
  \endlastfoot
  \textit{M. piperita} & 0 & 0 & 7.95 & 33.51 & 1.56 \\
   & 0 & 25 & 7.76 & 31.32 & 1.49 \\
   & 0 & 50 & 7.51 & 23.74 & 1.28 \\
   & 15 & 0 & 8.74 & 49.8 & 1.45 \\
   & 15 & 25 & 9.04 & 36.68 & 1.53 \\
   & 15 & 50 & 8.42 & 32.3 & 1.51 \\
   & 30 & 0 & 8.42 & 52.04 & 1.67 \\
   & 30 & 25 & 9.65 & 42.58 & 1.61 \\
   & 30 & 50 & 8.72 & 38.76 & 1.42 \\
  \midrule
  \textit{M. crispata} & 0 & 0 & 5.69 & 7.1 & 1 \\
   & 0 & 25 & 7.69 & 8.79 & 1.99 \\
   & 0 & 50 & 8.71 & 11.21 & 1.87 \\
   & 15 & 0 & 6.5 & 22.32 & 1.29 \\
   & 15 & 25 & 6.91 & 22.9 & 1.26 \\
   & 15 & 50 & 6.2 & 18.32 & 1.17 \\
   & 30 & 0 & 6.83 & 24.32 & 1.2 \\
   & 30 & 25 & 6.79 & 24.3 & 1.2 \\
   & 30 & 50 & 6.68 & 21.23 & 1.24 \\
  \midrule
  \textit{M. spicata} `Moroccan' & 0 & 0 & 6.65 & 7.61 & 0.99 \\
   & 0 & 25 & 5.95 & 6.61 & 0.99 \\
   & 0 & 50 & 8.09 & 8.1 & 1.05 \\
   & 15 & 0 & 7.67 & 50.91 & 1.27 \\
   & 15 & 25 & 7.08 & 40.62 & 1.2 \\
   & 15 & 50 & 7.51 & 32.49 & 1.12 \\
   & 30 & 0 & 8.07 & 56.92 & 1.16 \\
   & 30 & 25 & 7.51 & 41.76 & 1.17 \\
   & 30 & 50 & 8.01 & 42.05 & 1.22 \\
\end{longtable}
\end{singlespace}

\section{Interaction between mint species, growing media, and salt stress}

The effect of salinity depended on the growing medium and on the taxon. Because the taxa differed greatly in their absolute size, the interaction is most clearly seen in the proportional loss of plant height caused by the highest salinity, that is, the difference between the plants at 0 and 50~mM expressed as a percentage of the plants at 0~mM within the same growing medium. These proportional losses are given in Table~\ref{tab:decline} and shown in Figure~\ref{fig:decline}.

\begin{table}[htbp]
  \centering
  \caption[Proportional loss of plant height at 50~mM NaCl]{Loss of plant height at 50~mM NaCl relative to 0~mM within the same growing medium, as a percentage of the height at 0~mM.}
  \label{tab:decline}
  \begin{tabular}{l S[table-format=2.1] S[table-format=2.1] S[table-format=2.1]}
    \toprule
    & \multicolumn{3}{c}{Loss of plant height (\%)} \\
    \cmidrule(lr){2-4}
    Species & {Unamended} & {15~g biochar} & {30~g biochar} \\
    \midrule
    \textit{M. piperita}             & 30.7 & 23.3 & 18.0 \\
    \textit{M. crispata}             & 44.2 & 31.2 & 27.7 \\
    \textit{M. spicata} `Moroccan'   & 50.8 & 35.8 & 28.4 \\
    \bottomrule
  \end{tabular}
\end{table}

\begin{figure}[htbp]
  \centering
  \begin{tikzpicture}
    \begin{axis}[
      ybar,
      width=0.9\textwidth,
      height=7cm,
      bar width=14pt,
      ymin=0, ymax=60,
      ylabel={Loss of plant height at 50~mM (\%)},
      symbolic x coords={{\textit{M. piperita}},{\textit{M. crispata}},{`Moroccan'}},
      xtick=data,
      enlarge x limits=0.3,
      legend style={at={(0.5,-0.18)},anchor=north,legend columns=3,draw=none},
      ymajorgrids=true,
      grid style=dashed,
      nodes near coords,
      every node near coord/.append style={font=\footnotesize},
    ]
      \addplot[fill=gray!25] coordinates {({\textit{M. piperita}},30.7) ({\textit{M. crispata}},44.2) ({`Moroccan'},50.8)};
      \addplot[fill=gray!55] coordinates {({\textit{M. piperita}},23.3) ({\textit{M. crispata}},31.2) ({`Moroccan'},35.8)};
      \addplot[fill=gray!85] coordinates {({\textit{M. piperita}},18.0) ({\textit{M. crispata}},27.7) ({`Moroccan'},28.4)};
      \legend{Unamended peat, 15~g biochar, 30~g biochar}
    \end{axis}
  \end{tikzpicture}
  \caption[Loss of plant height at 50~mM NaCl by growing medium]{Loss of plant height caused by 50~mM NaCl, relative to the unsalted plants of the same growing medium, for each taxon and biochar rate.}
  \label{fig:decline}
\end{figure}

The pattern is consistent across the three taxa. In \textit{M. piperita}, salinity of 50~mM reduced plant height by 30.7\% in unamended peat, by 23.3\% with 15~g of biochar and by 18.0\% with 30~g, so that the loss with the higher biochar rate was about four tenths smaller than in unamended peat. In \textit{M. crispata} the corresponding losses were 44.2\%, 31.2\% and 27.7\%, and in \textit{M. spicata} `Moroccan' they were 50.8\%, 35.8\% and 28.4\%. In every taxon the proportional loss fell as the biochar rate rose, and in every taxon the 30~g rate gave the smallest loss.

The size of the benefit differed between the taxa. The reduction in the proportional loss between unamended peat and 30~g of biochar was 12.7 percentage points in \textit{M. piperita}, 16.5 in \textit{M. crispata} and 22.4 in \textit{M. spicata} `Moroccan'. The taxa that suffered most from salinity in unamended peat therefore gained most from biochar, and `Moroccan', which lost more than half of its height at 50~mM in unamended peat, lost only a little over a quarter at the higher biochar rate.

A second feature of the interaction is that biochar did not merely shift growth upwards by a constant amount. Had it done so, the proportional loss would have been similar in all three growing media. Instead the loss itself became smaller as the biochar rate rose, which indicates that the plants grown with biochar were less sensitive to the salt rather than simply larger to begin with. The same direction was seen in the chlorophyll content and, less consistently, in leaf length and width, whereas the Fv/Fm ratio was only slightly affected by either factor and showed no comparable pattern.

It should be noted that these proportional losses are calculated from single pots and are therefore descriptive. They are reported because they express the relationship between the two factors more directly than the absolute values, not because the differences between them have been shown to be statistically reliable.

\section{Statistical analysis of the effects of biochar and salinity}\label{sec:anova}

The effects of the biochar rate and of the salinity level were tested within each species by two-way analysis of variance on the pot means (Table~\ref{tab:anova}). Because there was one pot per treatment combination, the interaction between biochar and salinity could not be tested separately and these tests are exploratory.

Plant height was significantly affected by both factors in all three taxa ($p<0.01$). According to Tukey's test, plant height differed between all three salinity levels in every taxon, and between all biochar rates except between 15 and 30~g in \textit{M. spicata} `Moroccan' ($p=0.066$). The Fv/Fm ratio was significantly reduced by salinity in all taxa ($p<0.05$), but not significantly affected by the biochar rate ($p\geq0.10$). The chlorophyll content was significantly affected by both factors in \textit{M. piperita} (biochar $p=0.035$, NaCl $p=0.021$) and \textit{M. spicata} `Moroccan' (biochar $p=0.008$, NaCl $p=0.002$), but not in \textit{M. crispata} ($p=0.759$ and $p=0.144$), probably because of the unusual value of the unamended pot at 25~mM. Leaf width was significantly affected by both factors in \textit{M. piperita} and \textit{M. spicata} `Moroccan', and leaf length in \textit{M. spicata} `Moroccan' ($p<0.001$ for both); in \textit{M. crispata}, leaf length was significantly affected by salinity ($p=0.032$) but not by biochar ($p=0.052$). The relative water content was not included in the analysis of variance, for the reasons given above.

\begin{table}[htbp]
  \centering
  \small
  \caption[Two-way ANOVA of biochar rate and salinity]{Two-way analysis of variance (without interaction; 2 and 4 degrees of freedom) of the effects of biochar rate and salinity level on the pot means within each taxon. $F$ is the test statistic and $p$ the probability; values of $p<0.05$ were considered significant.}
  \label{tab:anova}
  \begin{tabular}{ll S[table-format=3.2,round-mode=places,round-precision=2] c S[table-format=3.2,round-mode=places,round-precision=2] c}
    \toprule
    Species & Variable & {$F$ (biochar)} & {$p$ (biochar)} & {$F$ (NaCl)} & {$p$ (NaCl)} \\
    \midrule
    \textit{M. piperita} & Plant height & 89.40 & $<$0.001 & 78.97 & $<$0.001 \\
     & Fv/Fm & 4.18 & 0.105 & 36.96 & 0.003 \\
     & Leaf length & 4.94 & 0.083 & 5.70 & 0.067 \\
     & Leaf width & 19.31 & 0.009 & 18.45 & 0.010 \\
     & Chlorophyll content & 8.65 & 0.035 & 11.81 & 0.021 \\
    \midrule
    \textit{M. crispata} & Plant height & 52.08 & 0.001 & 97.57 & $<$0.001 \\
     & Fv/Fm & 4.00 & 0.111 & 15.30 & 0.013 \\
     & Leaf length & 6.79 & 0.052 & 9.15 & 0.032 \\
     & Leaf width & 4.77 & 0.087 & 6.53 & 0.055 \\
     & Chlorophyll content & 0.30 & 0.759 & 3.27 & 0.144 \\
    \midrule
    \textit{M. spicata} `Moroccan' & Plant height & 34.41 & 0.003 & 52.15 & 0.001 \\
     & Fv/Fm & 1.72 & 0.288 & 15.13 & 0.014 \\
     & Leaf length & 63.66 & $<$0.001 & 115.07 & $<$0.001 \\
     & Leaf width & 13.85 & 0.016 & 43.09 & 0.002 \\
     & Chlorophyll content & 20.70 & 0.008 & 40.13 & 0.002 \\
    \bottomrule
  \end{tabular}
\end{table}
